\documentclass[11pt]{beamer}
\usepackage[utf8]{inputenc}
\usepackage{amsmath, amssymb, graphicx, listings, xcolor}

% Presentation Theme
\usetheme{Madrid}
\usecolortheme{default}

% Code Block Styling
\lstset{
    language=Python,
    basicstyle=\ttfamily\scriptsize,
    keywordstyle=\color{blue}\bfseries,
    commentstyle=\color{green!50!black},
    stringstyle=\color{red},
    showstringspaces=false,
    frame=single,
    backgroundcolor=\color{gray!10}
}

\title[Porkchop Plots \& NASA SPICE]{Astrodynamics 2026-1}
\subtitle{Characteristic Energy ($C_3$), Porkchop Plots, \& Computational Mission Design}
\author{Prof. Juan}
\institute{Universidad de Antioquia - Aerospace Engineering}
\date{Spring 2026}

\begin{document}

% --- Title Slide ---
\begin{frame}
    \titlepage
\end{frame}

% ==========================================
% SECTION 1: C3 ENERGY
% ==========================================
\section{Characteristic Energy ($C_3$)}

% --- Slide 1 ---
\begin{frame}{What is $C_3$?}
    When mission planners talk about launching to other planets, they rarely specify the initial launch velocity. Instead, they specify the \textbf{Characteristic Energy ($C_3$)}.
    
    \vspace{0.3cm}
    \textbf{Definition:} $C_3$ is the square of the Hyperbolic Excess Velocity ($v_\infty$).
    \begin{equation}
        C_3 = v_\infty^2
    \end{equation}
    
    \textbf{Physical Meaning:}
    \begin{itemize}
        \item It is a measure of the \textit{specific energy} required to escape Earth's gravity well and still have enough kinetic energy left over to reach the target planet.
        \item Units: $\text{km}^2/\text{s}^2$.
        \item If $C_3 = 0$, the spacecraft barely escapes Earth ($v_\infty = 0$), drifting into a solar orbit identical to Earth's.
        \item If $C_3 > 0$, the spacecraft escapes with residual speed to travel outward to Mars or inward to Venus.
    \end{itemize}
\end{frame}

% --- Slide 2 ---
\begin{frame}{Why Use $C_3$ Instead of $\Delta V$?}
    Why do we feed $C_3$ into our numerical solvers instead of just calculating $\Delta V$?
    
    \vspace{0.3cm}
    \textbf{1. It isolates the interplanetary problem:} \\
    $\Delta V$ depends on the altitude of your specific parking orbit ($r_p$). $C_3$ is an absolute energy state independent of the launch vehicle's parking orbit. 
    
    \vspace{0.3cm}
    \textbf{2. It dictates Launch Vehicle Selection:} \\
    Companies like SpaceX and ULA publish performance curves for their rockets based strictly on $C_3$. If our mission requires a $C_3$ of $15 \text{ km}^2/\text{s}^2$, we look at the rocket's manual to see exactly how much payload mass it can push to that energy level.
\end{frame}

% ==========================================
% SECTION 2: THE PORKCHOP PLOT
% ==========================================
\section{The Porkchop Plot}

% --- Slide 3 ---
\begin{frame}{The Reality of Interplanetary Flight}
    In Phase 1 of Patched Conics, we assumed planets move in perfect, coplanar circles. \textbf{This is never true in reality.}
    
    \vspace{0.3cm}
    Because target planets have eccentric and inclined orbits:
    \begin{itemize}
        \item A perfect 180-degree Hohmann transfer almost never exists.
        \item The energy ($C_3$) required to reach the target changes drastically depending on exactly \textit{which day} we leave and \textit{which day} we arrive.
    \end{itemize}
    
    \vspace{0.3cm}
    \textbf{The Solution:} We use a computer to solve Lambert's Problem for \textit{thousands} of different departure and arrival date combinations to find the cheapest path.
\end{frame}

% --- Slide 4 ---
\begin{frame}{Mapping the Energy Landscape}
    When we plot the $C_3$ energy required for all these thousands of date combinations on a 2D contour map, a distinct shape emerges: \textbf{The Porkchop Plot}.
    
    \vspace{0.3cm}
    \textbf{How to read it:}
    \begin{itemize}
        \item \textbf{X-Axis:} Departure Date (from Earth).
        \item \textbf{Y-Axis:} Arrival Date (at target planet).
        \item \textbf{Contour Lines:} The $C_3$ energy required for that specific flight.
        \item \textbf{The "Valleys":} The dark blue centers of the porkchop shape represent the optimal Launch Windows—the dates requiring the absolute minimum fuel.
    \end{itemize}
\end{frame}

% ==========================================
% SECTION 3: THE PYTHON / SPICE CODE
% ==========================================
\section{Computational Architecture}

% --- Slide 5 ---
\begin{frame}[fragile]{The Code: 1. SPICE Kernels \& Time Grid}
    To build our Porkchop Plot, we first tell Python which dates to check. We create arrays of time stepping forward by 2 days.
    
\begin{lstlisting}
# Define search windows for Departure and Arrival
dep_start = spy.str2et('2026-09-01')
dep_end = spy.str2et('2026-11-01')

arr_start = spy.str2et('2027-05-01')
arr_end = spy.str2et('2027-08-01')

# Create arrays stepping by 86400 * 2 seconds
departure_dates = np.arange(dep_start, dep_end, step_size)
arrival_dates = np.arange(arr_start, arr_end, step_size)
\end{lstlisting}
    \textit{Note: \texttt{str2et} converts Gregorian dates into Ephemeris Time (seconds past J2000), which SPICE requires for math.}
\end{frame}

% --- Slide 6 ---
\begin{frame}[fragile]{The Code: 2. Fetching True Ephemeris}
    Inside our nested loops, we grab the true 3D state vectors for the planets on those specific dates using \texttt{spkgeo}.
    
\begin{lstlisting}
for t_dep in departure_dates:
    # Get Earth's true vector at departure (R1)
    state_earth, _ = spy.spkgeo(399, t_dep, 'ECLIPJ2000', 10)
    r1 = state_earth[:3]
    v_earth = state_earth[3:]
    
    for t_arr in arrival_dates:
        # Get Target's true vector at arrival (R2)
        state_target, _ = spy.spkgeo(TARGET_ID, t_arr, 'ECLIPJ2000', 10)
        r2 = state_target[:3]
\end{lstlisting}
    \textit{We now have $R_1$, $R_2$, and the Time of Flight ($t_{arr} - t_{dep}$). We are ready to solve Lambert's Problem!}
\end{frame}

% --- Slide 7 ---
\begin{frame}[fragile]{The Code: 3. Solving Lambert's Problem}
    We feed $R_1$, $R_2$, and Time of Flight into the \texttt{poliastro} Lambert solver. It returns the exact velocity we need ($v_{tx1}$) to hit the target.
    
\begin{lstlisting}
# poliastro solves the boundary value problem
solutions = list(izzo.lambert(MU_SUN, r1, r2, tof))

# Extract the velocity for the direct (0-rev) transfer
v_tx1, v_tx2 = solutions[0]
v_tx1 = v_tx1.value 
\end{lstlisting}
    \vspace{0.2cm}
    \textbf{The Physics:} The solver calculates the elliptical arc connecting the two planets and outputs the heliocentric velocity the spacecraft must have as it leaves Earth's Sphere of Influence.
\end{frame}

% --- Slide 8 ---
\begin{frame}[fragile]{The Code: 4. Calculating $C_3$ \& Optimizing}
    Finally, we calculate the $C_3$ for that specific date combination. If it is lower than our previous best, we save it.
    
\begin{lstlisting}
# 1. Hyperbolic Excess Velocity Vector
v_inf_vec = v_tx1 - v_earth

# 2. C3 is the magnitude of v_infinity squared (dot product)
c3 = np.dot(v_inf_vec, v_inf_vec)

# 3. Save if it's the lowest energy found so far
if c3 < best_c3:
    best_c3 = c3
    best_dep_time = t_dep
    best_arr_time = t_arr
\end{lstlisting}
    \vspace{0.2cm}
    \textit{By the end of the loop, Python has swept the entire Porkchop Plot and found the exact bottom of the energy valley!}
\end{frame}

% ==========================================
% SECTION 4: APPLICATION - VENUS
% ==========================================
\section{Application: The Venus Launch Window}

% --- Slide 9 ---
\begin{frame}[fragile]{Adapting the Code for Venus}
    To find the launch window for our Earth-to-Venus mission, we only need to make two changes to the script's initialization parameters:
    
    \vspace{0.3cm}
    \textbf{1. Change the Target ID:} \\
    Venus's NAIF ID is \texttt{299} (or the barycenter \texttt{2}).
\begin{lstlisting}
TARGET_ID = 299
\end{lstlisting}

    \vspace{0.3cm}
    \textbf{2. Set the Time Grid for Venus:} \\
    Based on our analytical synodic period calculations (584 days), the next major window to Venus opens in early 2028.
\begin{lstlisting}
dep_start = spy.str2et('2028-02-01')
dep_end = spy.str2et('2028-04-01')

arr_start = spy.str2et('2028-07-01')
arr_end = spy.str2et('2028-10-01')
\end{lstlisting}
\end{frame}

% --- Slide 10 ---
\begin{frame}{Closing the Loop: From Python to $\Delta V$}
    Suppose your Python script sweeps the 2028 Venus window and finds an optimal launch date requiring $C_3 = 6.25 \text{ km}^2/\text{s}^2$. 
    
    \vspace{0.2cm}
    How do we translate this software output back to the physical $\Delta V$ required for our spacecraft parked in a 300 km LEO ($r_p = 6678$ km)?
    
    \vspace{0.2cm}
    \textbf{1. Calculate Periapsis Velocity ($v_p$):}
    \begin{equation*}
        v_p = \sqrt{C_3 + \frac{2\mu_\oplus}{r_p}} = \sqrt{6.25 + \frac{2(398600)}{6678}} = \mathbf{11.21 \text{ km/s}}
    \end{equation*}
    
    \textbf{2. Calculate Engine Burn ($\Delta V$):}
    \begin{equation*}
        v_c = \sqrt{\frac{\mu_\oplus}{r_p}} = \mathbf{7.73 \text{ km/s}} \quad \Rightarrow \quad \Delta V = 11.21 - 7.73 = \mathbf{3.48 \text{ km/s}}
    \end{equation*}
    
    \textit{This mathematically bridges your numerical SPICE search with your analytical Patched Conic theory!}
\end{frame}

\end{document}